\lesson{2}{Mar 04 2022 Fri (19:03:22)}{Create QE from roots and LC}{Unit 4}

\begin{theorem}[Fundamental Theorem of Algebra]
  \label{thrm:fundamental_theorem_of_algebra}

  Any polynomial of degree $n > 0$ has at least one zero (real or non-real).
\end{theorem}

\begin{theorem}[Factor Theorem]
  \label{thrm:factor_theorem}

  If $r$ is a zero of the polynomial $P(x)$, then $x - r$ is a factor of $P(x)$.
  Conversely, if $x - r$ is a factor of  $P(x)$, then $r$ is a zero of $P(x)$.
\end{theorem}

\begin{example}
  \label{exm:factor_theorem_example}

  Consider the polynomial $P(x) = x^{2} + 3x + 2$.

  Because

  \begin{equation}
    \label{eqt:factor_theorem_example_1}
    \begin{split}
      P(-1) = (-1)^{2} + 3(-1) + 2 = 1 -3 + 2 = 0 \\
    \end{split}
  .\end{equation} 

  we have that $-1$ is a zero of the polynomial $P(x)$.

  Hence, the \textbf{Factor Theorem} (\ref{thrm:factor_theorem}) guarantees that
  $x - (-1) = x + 1$ is a factor of $P(x)$.

  That is, the factor theorem guarantees that there exists a polynomial $Q(x)$
  such that

  \begin{equation}
    \label{eqt:factor_theorem_example_2}
    \begin{split}
      P(x) = (x + 1)Q(x)
    \end{split}
  .\end{equation}

  That can be confirmed directly. Since $P(x) = x^{2} + 3x + 2 = (x + 1)(x + 2)$,
  we have that $Q(x) = x + 2$
\end{example}

\begin{remark}
  \label{rmk:factor_theorem_remark}

  Given a polynomial
  $P(x) = a_n x^{n} + a_{n - 1} x^{n - 1} + \ldots + a_1 x + a_0$ of degree
  $n > 0$, we can apply the \textbf{Fundamental Theorem of Algebra}
  (\ref{thrm:fundamental_theorem_of_algebra}) and facto theorem $n$ 
  times to obtain what is commonly referred to as the complete factorization of
  $P(x)$. Namely, we can write $P(x)$ as follows:

  \begin{equation}
    \label{eqt:factor_theorem_remark_1}
    \begin{split}
      P(x) = a_n (x - r_1)(x - r_2) \cdots (x - r_n)
    \end{split}
  .\end{equation}

  where $a_n$ is the leading coefficient of $P(x)$ and $r_1, r_2, \ldots r_n$ 
  are the $n$ zeroes of $P(x)$. This fact is sometimes referred to as the
  \textbf{Linear Factors Theorem}
\end{remark}

Let's write a quadratic equation whose roots are $-1 \quad -3$ and whose leading
coefficient is $2$.

We use the \textbf{Factor Theorem} (\ref{thrm:factor_theorem}), which states
that if $k$ is a root of the polynomial equation $P(x) = 0$, then $x - k$ is a
factor of the polynomial $P(x)$.

In our problem, we have $P(-1) = 0 \quad P(-3) = 0$. Thus, $x + 1 \quad x + 3$
are both factors of the quadratic polynomial. Since the coefficient is $2$, we
can write the polynomial as:

\begin{equation}
  \label{eqt:creating_polynomial_from_roots}
  \begin{split}
    2(x + 1)(x + 3) = 2x^{2} + 8x + 6
  \end{split}
.\end{equation}

\newpage
